% !PRINTED 209
\weberpage{209}

That conversely every such expression, if it is homogeneous in
the coefficients $a$, represents an invariant with integer coefficients,
is clear of itself\footnote{The theory of invariants is presented in detail in
the works: \gesperrt{Clebsch}, ``Theorie der bin\"aren algebraischen Formen''.
Leipzig 1872. \gesperrt{Fa\`a di Bruno}, ``Einleitung in die Theorie der bin\"aren
Formen''. German by \gesperrt{Walter}. Leipzig 1881. \gesperrt{P. Gordans}, ``Vorlesungen
\"uber Invariantentheorie'', edited by \gesperrt{Kerschensteiner}. Leipzig
1885. Cf. also \gesperrt{Franz Meyer}, ``Bericht \"uber den gegenw\"artigen Stand
der Invariantentheorie'' in the Jahresbericht der deutschen Mathematiker-
Vereinigung 1890/91 (Berlin 1892).}.

% Printer's furniture excluded by policy: footer "Weber, Algebra. I." and gathering number 14.
